\scanpage{85}
% loc 1–5
et on vérifie l'associativité. De plus, on prolonge $\operatorname{Tr}$ à $V$ par linéarité, et $\det$, \deleted{par} d'abord sur $\underbrace{k[\underline u]}_{A_1}\subset M$ et $\underbrace{k[\underline v]}_{A_2}\subset M$ (où c'est \deleted{$\det(u'$} la norme des algèbres quadratiques, $\operatorname{Tr}$ y étant la Trace) puis sur
% loc 6–7
\[ M=k[u]\oplus k[v]\qquad\text{(}\inserted{\text{iso}}\ \text{du }\text{pt de vue}\ \text{additif)} \]
% loc 8
par
% loc 9
\[ \det(u_1^{\deleted{\unsure{2}}}+v_1)=\underbrace{\det u_1}_{N_{A_1/k}(u_1)}+\underbrace{\det v_1}_{N_{A_2/k}(v_1)}+\operatorname{Tr}u_1\operatorname{Tr}v_1-\operatorname{Tr}(u_1v_1) \tag*{(10)} \]
% loc 10–13
Ceci dit, ``on \deleted{voit}\inserted{vérifie} que'' $\det$ est quadratique et multiplicatif, que l'on a la relation de Hamilton-Cayley et la relation (3).

% loc 14–18
En fait, je n'ai pas eu à faire la vérification, car \deleted{\unsure{je sai|je sais}} il suffit de la faire (principe de prolongement des id alg) \deleted{pour} en supposant
% loc 19
\[ \delta\overset{\text{dfn}}{=}b^2-4c\ \text{inv.},\quad \delta'\overset{\text{dfn}}{=}b'^2-4c'\ \text{inv.} \]
% loc 20–22
(i.e. $A_1,A_2$ étales sur $k$), et on \deleted{sait} voit que, \deleted{pa}\inserted{pour $b,c,b',c'$ donnés ainsi,} l'on \deleted{peut} réaliser \deleted{(pour alors (pour pour) le \unsure{système|septime|répartition} $(b,c$}
% archive sheet number: 84
